\section{Discussion}
\label{sec:discussion}

\paragraph{Why the collapse is generic.} Two source-free methods, four non-adaptive random forests and two test rigs lose
between 22 and 41 points when the target bearings are disjoint from the source bearings. The probe of
Section~\ref{sec:res-probe} explains why a specific algorithm is not the culprit: the source features identify individual
specimens within a class almost perfectly, so any classifier trained on them can separate the source classes by identity.
With three bearings per class, identity is not a weak shortcut competing with fault physics --- it is a perfect one. Nothing
in an unsupervised adaptation objective penalises it, because adaptation refines the source model's own decisions rather
than questioning what they are based on. This also predicts the shape of the failure we observe: not degraded per-window
accuracy, but a confident label applied to a whole unseen specimen.

\paragraph{Why physics gating helps, and why only this much.} The gates that are safe on healthy machinery deliver 4--7
points, in every seed, in 9 of 10 splits, and essentially never cost anything. That combination --- small, consistent,
almost risk-free --- is what a correct but partial prior looks like. The ceiling analysis quantifies the ``partial''. First,
a pseudo-label filter can only repair the filter-recoverable part of the collapse, a median \SI{82}{\percent} on Paderborn
but as little as \SI{48}{\percent} on some splits. Second, a kinematic gate acts only where fault lines are actually
detectable, which on this dataset is about a quarter of the faulty bearings. The product of those two factors is a small
number, and no amount of tuning changes it. Reporting a gain without the oracle gap beside it hides exactly this.

\paragraph{Safety is a property to measure, not to assume.} The raw-spectrum band rule accepts \SI{45}{\percent} of healthy
recordings as faulty, and inside the adaptation loop it makes accuracy worse. The envelope rules accept at most
\SI{2}{\percent} and help. The ordering of the two effects is the same as the ordering of the false-acceptance rates, and
both are measurable before any adaptation experiment is run, on healthy recordings from the target rig. A physics module
introduced into an adaptation loop should be characterised this way first; end accuracy on a leaky protocol cannot
distinguish a safe prior from an unsafe one.

\paragraph{A shallow baseline belongs in every such comparison.} On bearing-disjoint targets a random forest on ten
time-domain statistics --- no adaptation, no target data, no deep network --- exceeds the adapted source-free method by a
median of 8.9 points across the ten splits, and matches the gated method. We do not present this as a better diagnostic
method; it collapses for the same reason and labels bearings as blocks in the same way. It is a calibration point. Any
claim that an adaptation procedure has learned transferable fault structure should clear the bar set by features that
encode no transfer mechanism at all.

\paragraph{Is this leakage, or simply a different task?}
A reader may object that adapting to the same physical bearings at a new operating condition is a legitimate task: the
machine is the same, only its duty point has changed, and that is a real deployment scenario for a permanently instrumented
asset. We agree, and the objection sharpens rather than weakens the argument. Two scenarios are conflated in the current
literature. In the first, the monitored bearing is present in training, and same-bearing evaluation estimates performance
honestly. In the second --- diagnosing a replacement bearing, a sister machine, or a fleet --- the specimen at test time was
never seen, and same-bearing evaluation over-estimates performance by 22--41 points on the rigs measured here. The failure is
therefore not that the published protocol is wrong in itself, but that the reported number is read as evidence for the
second scenario while being measured under the first. We use ``leakage'' in the sense of Kapoor and Narayanan
\cite{kapoor2023leakage}: information shared between training and evaluation that inflates the estimate relative to the
deployment the claim addresses. What the adapted model exploits is a shortcut in exactly the sense of Geirhos et al.
\cite{geirhos2020shortcut} --- a rule that succeeds on the benchmark and fails under the intended shift.

\paragraph{Threats to validity.} Three deserve naming. The clean and leaky targets differ in their bearing sets, so a split
where the clean bearings happen to be harder inflates $\Delta_{\mathrm{leak}}$; running both directions on the folds and ten
random splits addresses this, and the split-level spread is reported rather than hidden. The two Paderborn folds are not
perfectly symmetric (fold B's source has two outer-race bearings rather than three), which is visible in the decomposition
and is why the folds are reported separately from the ten random splits. Finally, the HUST replication changes bearing size
along with identity; it strengthens the leakage conclusion but cannot isolate identity alone, and we do not use it for any
claim about gating.
